%! TEX root = ../main.tex

\chapter{SSD Online Calibration}
\label{chap:further-improvements}

The previous sections in this work have demonstrated the excellent capabilities
of the \UMD to predict the primary type of a \CR by accessing the muonic 
component of \EASs directly. The presented discrimination framework is able to 
identify \todo{SE} of all photon simulations without falsely tagging any 
proton simulations as candidates at energies \todo{energy}. 

With these performance figures in mind, the metric to be optimized next becomes 
the size of the search sample. This can be increased in two main ways. First, 
the trigger efficiency of the \SD can be improved, resulting in more observed 
\EASs, especially at lower energies. Second, the detector area can be increased
via extending the analysis e.g.~to the \SD-1500 by virtue of performing primary 
discrimination using a joint analysis of the \WCD and \SSD signal in an event.

The foundation needed for both cases is the implementation of the AugerPrime 
detectors in the event detection chain. For the \SSD in particular, this
requires an algorithm that puts measured signals into a physical context. Such
an algorithm has been presented in \cref{sssec:ssd}, but has some shortcomings.
Notably, it is not a in-situ algorithm, limiting its' temporal resolution to an
order of hours. This impedes an accurate calculation of the exposure of the 
\SSD-array, and constrains the effectiveness of standalone event triggers 
relying on \SSD data, both critical in extending the photon search to the 
\SD-1500 array.

In the following, we address these issues by implementing an 
online\footnote{Meaning in-situ, as opposed to histogram-fitting, which happens
offline (ex-situ).} calibration algorithm for the \SSD-array. We present the 
working procedure of the routine in \cref{sec:ssd-online-calibration-algorithm},
and state the preliminary performance in 
\cref{sec:ssd-online-calibration-performance}. Finally, we present a 
minute-accurate monitoring framework of the \SSD-array in 
\cref{sec:ssd-online-calibration-monitoring}, which relies on the data produced
by the online calibration alogrithm.

disclaimer, citing gap notes, yada yada

\section{Calibration Algorithm}
\label{sec:ssd-online-calibration-algorithm}

\subsection{Algorithm Routine}

We adopt the idea of a rate-based algorithm from the already implemented \WCD
online calibration (Ref.~\cref{sssec:wcd}). In this spirit, we implement a 
single-bin calibration trigger with a trigger threshold $T_{70}$ whose trigger
rate is monitored for an integration window $t_\mathrm{cal.}$. After 
$t_\mathrm{cal.}$ has passed, the trigger threshold is adjusted by a value of
$\delta$ such that the observed trigger rate, $f_{70}$ converges to a target 
rate $f_{70}^\mathrm{target}$. For compatibility reasons, we chose 
$f_{70}^\mathrm{target} = \SI{70}{\hertz}$. In order to increase convergence 
speed, we also adjust the magnitude of $\delta$ based on the measured difference
$\Delta r = \left|f_{70} - f^\mathrm{target}_{70}\right|$. The control diagram
for the process is shown in \cref{fig:control-diagram}.

\begin{figure}
    \centering
    
\vspace{0.5cm}
\begin{tikzpicture}[font=\footnotesize]

% Start block
\node[draw,
    rounded corners,
    rectangle,
    align=center,
    text width=2.1cm,
    minimum width=2.2cm,
    minimum height=1cm] (block0) 
    {
    Adjust by \\
    $t_\mathrm{cal.}+\SI{5}{\second}$ \\
    $\delta-0.01\,\mathrm{ADC}$
    };

\node[draw,
    align=center,
    text width=4.2cm,
    above=of block0,
    minimum width=3cm,
    minimum height=1cm
    ] (block1) {
    $\Delta r > 0:\,T_{70}^\mathrm{new} = T_{70}^\mathrm{old} + \delta$
    $\Delta r < 0:\,T_{70}^\mathrm{new} = T_{70}^\mathrm{old} - \delta$
    };

\node[draw,
    shape=diamond,
    left=0.6cm of block1,
    minimum height=1cm
    ] (block3) { 
    Wait $t_\mathrm{cal.}$
    };

\node[draw,
    align=center,
    minimum height=1cm,
    below=of block0
    ] (block4) {
    $\Delta r < \sigma_r$
    };

\node[draw,
    align=center,
    left=1.54cm of block4,
    minimum width=3cm,
    minimum height=1cm
] (block5) {
	$\Delta r = \left|f_{70} - \SI{70}{\hertz}\right|$ \\ $\sigma_{r} = \sqrt{n\,/\,t_\mathrm{cal.}}$
    };

\node[draw,
    rounded corners,
    rectangle,
    right=0.4cm of block0,
    text width=2.1cm,
    align=center,
    minimum width=2.2cm,
    minimum height=1cm] (block6) 
    {
    Set increment\\
$\delta = 0.1\,\mathrm{ADC}.$
    };

\node[draw,
    align=center,
    minimum height=1cm,
    below=1.017cm of block6
    ] (block7) {
    $\Delta r < \SI{5}{\hertz}$
    };

\node[draw,
    rounded corners,
    rectangle,
    right=0.4cm of block6,
    text width=2.1cm,
    align=center,
    minimum height=1cm] (block8) 
    {
    Leave $t_\mathrm{cal.}$ and $\delta$ invariant
    };

\node[draw,
    align=center,
    minimum height=1cm,
    below=0.968cm of block8
    ] (block9) {
    $\Delta r < \SI{20}{\hertz}$
    };

\node[draw,
    rounded corners,
    rectangle,
    right=0.4cm of block8,
    text width=2.1cm,
    align=center,
    minimum height=1cm] (block10) 
    {
    Set values \\
    $t_\mathrm{cal.} = \SI{10}{\second}$, \\
    $\delta = 1.0\,\mathrm{ADC}$.
    };

\node[draw,
    rounded rectangle,
    above=3.55cm of block8,
    align=center,
    minimum width=2.2cm,
    minimum height=1cm] (block13) 
    {START};

\node[draw,
    rounded rectangle,
    left=1.9cm of block13,
    align=center,
    minimum width=2.2cm,
    minimum height=1cm] (block2) 
    {
    Initialize $T_{70} = 50\,\mathrm{ADC}$, $t_\mathrm{cal.} = \SI{10}{\second}$, and $\delta = 0.2\,\mathrm{ADC}.$
    };

\node[coordinate,right=0.54cm of block1] (block11) {};
\node[coordinate,right=2.764cm of block11] (block12) {};

% Arrows
\draw[-latex] (block2) edge (block3);
\draw[-latex] (block3) edge (block5);
\draw[-latex] (block5) edge (block4);
\draw[-latex] (block13) edge (block2);

\draw[-latex] (block4) edge node[pos=0.45,fill=white,inner sep=2pt]{Yes}(block0);
\draw[-latex] (block4) edge node[pos=0.45,fill=white,inner sep=2pt]{x}(block7);
\draw[-latex] (block7) edge node[pos=0.4,fill=white,inner sep=2pt]{Yes}(block6);
\draw[-latex] (block7) edge node[pos=0.45,fill=white,inner sep=2pt]{x}(block9);
\draw[-latex] (block9) edge node[pos=0.4,fill=white,inner sep=2pt]{Yes}(block8);
    
\draw[-latex] (block9) -| (block10)
    node[pos=0.15,fill=white,inner sep=1]{No};

\draw[-latex] (block1) edge (block3);
\draw[-latex] (block0) edge (block1);
\draw (block8) edge (block12);
\draw (block6) -- (block11);
\draw[-latex] (block10) |- (block1);

\end{tikzpicture}
\vspace{0.5cm}

    \caption{Control diagram for the rate-based SSD online calibration algorithm}
    \label{fig:control-diagram}
\end{figure}


For example, if $730$ triggered events are registered in the first 
$t_\mathrm{cal.} = \SI{10}{\second}$ after initialization, we calculate 
$\Delta f$ to be \SI{3}{\hertz}. This exceeds $\sigma_r\approx\SI{2.7}{\hertz}$,
but not \SI{5}{\hertz}. We therefore set $\delta=\SI{0.1}{ADC}$, and, because 
$\Delta f > 0$, increase $T_{70}$ by $\delta$. We then collect the number of 
events that satisfy the updated trigger threshold $T_{70} = \SI{50.1}{ADC}$ for
the following integration window ($t_\mathrm{cal.}=\SI{10}{\second}$), and 
repeat the above procedure.   

In this way, the trigger threshold $T_{70}$ is adjusted until a \SI{70}{\hertz} 
event rate of the single-bin trigger is reached. Once this is the case, the 
integration interval $t_\mathrm{cal.}$ is increased in 
\SI{5}{\second}-increments to minimize the influence of Poissonian fluctuations 
on the end result. In order for the algorithm to run stable under varying
conditions we clip the calibration parameters 
$\delta / \mathrm{ADC} \in [0.01, 1.00]$ and 
$t_\mathrm{cal.} / \mathrm{s} \in [10, 61]$ respectively. Whenever the algorithm
completes a maximum integration interval of $\SI{61}{\second}$, the current 
value of $T_{70}$ is exported to other processes of the \UUB.

\subsection{Feasibility Tests}

Assuming constant electronic gain and atmospheric conditions, this algorithm is 
guaranteed to produce a stable value of $T_{70}$, which will in turn relate
to the value of \Imip by a constant factor

\begin{equation}
\label{eq:mip-online-calibration-factor}
\Imip\approx \MipOnline = T_{70}\,\big/\,C.
\end{equation}

However, temperature and pressure changes in the air complicate this picture. 
The attenuation of the \EM component of \EASs varies with a denser atmosphere, 
and the rate of electrons impinging on the \SSD changes. This results in a 
complicated dependency of $C$ in \cref{eq:mip-online-calibration-factor} on 
temperature, pressure, humidity, and air composition, none of which are 
monitored on a station-level. 

This imperfect knowledge decreases the resolution of the rate-based estimator 
\MipOnline, and may limit the ability of the \SSD to participate in the event 
detection chain. To ensure a stable operation, we desire an estimator that is 
precise. It need not be very accurate, as biases can be corrected 
for.\footnote{As is done e.g.~when converting from omnidirectional to 
vertical muon peak (\cref{sssec:reconstruction-histogram-fitter}).}
This is to say that the minute-to-minute fluctuation in \MipOnline should be 
minimal, whereas the bias of \MipOnline is not too important as long as it is 
consistent across stations.

Three main factors drive $\upsigma(\MipOnline)$; (a) the aforementioned 
non-constancy of $C$ due to changing atmospheric conditions; (b) intra- and 
inter-station fluctuations; and (c) the Poissonian sampling error on the number 
of calibration triggers. 

In order to verify whether these systematics can be controlled to a suitable 
degree, we perform feasibility tests using $870,\!000$ \SSD muon histograms 
(\cref{sssec:muon-histograms}) acquired from \TTHREE-events of $1,\!311$ 
distinct \SD-stations in the months June and December 2023. 

\subsubsection{Determination of the ``True'' MIP Peak}

On an event basis, we first draw sample muon events from the original SSD 
histogram to obtain a finer binning. Doing this, we neglect events from the 
first 12 bins ($\leq24$\,ADC). We assume a uniform distribution of events per 
original SSD bin and conserve counts (and thus rates) of the original histogram.
Next, we fit a parabola from $-6$\,ADC to $+15$\,ADC relative to the peak 
location of the resampled histogram and calculate the ordinate of the maximum of
the recovered parabola. This is qualitatively equivalent to the \Offline 
strategy (\cref{sssec:reconstruction-histogram-fitter}) to obtain the value of 
\MipOffline. Technically, resampling is not yet required, since fitting the 
original histogram must, on average, result in the same value for \MipOffline. 
Nevertheless, we perform the resampling anyways, as it will allow for a finer 
resolution in rate determination later on. In any case, resampling only changes 
the y-scale of the histogram, and does not introduce a bias on \MipOffline. An 
example event with the corresponding fits of \MipOffline are shown in 
\cref{fig:ssd-histograms}.

\begin{figure}
  \centering
  \def\w{0.48}
  \subfloat{\includegraphics[width=\w\textwidth]{ssd-calibration/true_mip.png}
  \label{fig:ssd-histograms}
  }\hfill
  \subfloat{\includegraphics[width=\w\textwidth]{ssd-calibration/rates.png}
  \label{fig:array-rates}
  }
  \caption[]{\subref{fig:ssd-histograms} An example \SSD \TTHREE histogram 
  (black) and its resampled equivalent (red). Fitting the histogram peak (dashed
  lines) of either histogram yields compatible results. \subref{fig:array-rates}
  The expected rate $r^+$ of events with maximum pulse height above a given 
  threshold. The statistical error from sampling fluctuations is Poissonian and 
  given as $\sqrt{r^+}$. The systematic error is the standard deviation of $r^+$
  over the entire dataset. It originates from minute-to-minute fluctuations in 
  $\Imip$ per station.}
  \label{fig:histograms-and-rate}
\end{figure}

\subsubsection{SSD-Rate/Threshold Relationship}

After determining the value of \MipOffline, we calculate the number of 
calibration events in the resampled SSD histogram that satisfy a threshold of 
$C = \{1.00, 1.05, \ldots, 5.00\}$ in units of $\MipOffline$. Dividing this by 
the acquisition time of \SI{61}{\second} yields an estimate of the calibration 
trigger rate $r^+$ at the time of the \TTHREE.

With this, we calculate the average rate $\langle r^+_i\rangle$ across all 
events for station $i$, and $\langle r^+\rangle$ across the entire array after
pruning events $j$ where 
$|\langle r^+_i\rangle-r^+_{i,j}| > 3\,\upsigma(r^+_i)$. The resulting 
rate/threshold relationship is shown in \cref{fig:array-rates}.

Next, we investigate the performance of the estimator \MipOnline gained from 
inverting the rate/threshold relationship above. In particular, we compare the
fitted value \MipOffline with the threshold read off from 
\cref{fig:array-rates} and converted via 
\cref{eq:mip-online-calibration-factor}.

We observe a positive correlation for all considered thresholds. The degree of 
correlation increases from $\rho\approx0.5$ to $\rho\approx0.75$ when increasing
the $C$ from $1$ to $3$ (or decreasing the rate from \SI{200}{\hertz} to 
\SI{10}{\hertz}). This is due to the smaller influence of low-energy processes 
skewing the histogram distribution at higher ADC counts. By selecting events 
with higher ionization power we mitigate effects of readout noise and physical
fluctuations. The fact that one can in principle estimate \MipOffline by 
building an integral rate is clear from this result.

We aim to quantize the smearing in the rate/threshold relationship by splitting 
the set of histograms into two distinct subsamples with maximum difference in 
temperature range. The first sample consists of \TTHREE events recorded between
10:00 to 14:00 local time in June 2023, whereas the second sample comprises 
events recorded between 22:00 to 02:00 during december nights.

\begin{figure}
  \centering
  \def\w{0.48}
  \subfloat{\includegraphics[width=\w\textwidth]{ssd-calibration/summer_days.png}
  \label{fig:summer-days}
  }\hfill
  \subfloat{\includegraphics[width=\w\textwidth]{ssd-calibration/winter_nights.png}
  \label{fig:winter-nights}
  }
  \caption[]{Bias and resolution of \MipOnline with respect to \MipOffline shown
  for events during \subref{fig:summer-days} summer days or 
\subref{fig:winter-nights} for a varying rate/threshold relationship.}
  \label{fig:drift-plots}
\end{figure}

The bias and resolution for both data sets are shown and \cref{fig:drift-plots}.
In either case, the 

% continue rewrite here

From \cref{fig:correlations} and \cref{fig:drift-plots}, it follows that (a) is smaller at higher calibration trigger thresholds and reads about 5 to 8\% at values above $\sim2\,\MipOffline$, (b) increases due to diminishing rates, contributing an error of ${>}10\%$ at $5\,\Imip$. At intermediate cuts of 1.5 to $3\,\Imip$ the combined error is calculated to be ${\sim}6\%$, which is sufficient for stable operation in the SD-1500 array. The behaviour of $\sigma_{\Delta\MipOnline}$ (which is the expected error on \MipOnline for all events of a single station) across all considered thresholds is shown alongside the systematic difference between \MipOffline and \MipOnline averaged over all stations in \autoref{fig:performance-plot}

\subsection{Integration Tests}

\section{Performance}
\label{sec:ssd-online-calibration-performance}

\subsection{Preprocessing}
\subsection{Station-to-Station Fluctuations}
\subsection{Temperature Fluctuations}
\subsection{Seasonal Fluctuations}
\subsection{Bias and Resolution}

\section{Monitoring}
\label{sec:ssd-online-calibration-monitoring}

\subsection{Data Set}
\subsection{SSD-array Uptime}
\subsection{Calibration Issues}
% put all the different shit like sunrise/sunset noise here

\section{Summary}

